% !TeX root = ../../progress-sem1.tex
% 1.x แผนการดำเนินงาน (Task breakdown + Gantt)
\section{แผนการดำเนินงาน}
\subsection{Task Breakdown}
\begin{enumerate}
  \item \textbf{Documentation}
    \begin{itemize}
      \item ศึกษาค้นคว้าหัวข้อโครงงาน
      \item จัดทำเอกสาร Project Idea
      \item เอกสารและงานวิจัยที่เกี่ยวข้อง
      \item จัดทำเอกสาร Proposal report
      \item นำเสนอ Proposal
      \item จัดทำเอกสาร Progress report ภาคเรียนที่ 1
      \item นำเสนอ Progress ภาคเรียนที่ 1
      \item จัดทำเอกสาร Progress report ภาคเรียนที่ 2
      \item นำเสนอ Progress ภาคเรียนที่ 2
      \item จัดทำเอกสาร Final Report
      \item นำเสนอ Final Defense
    \end{itemize}

  \item \textbf{Requirements Analysis \& System Design}
    \begin{itemize}
      \item สัมภาษณ์ผู้มีส่วนได้ส่วนเสีย
      \item ระบุ Requirements
      \item ระบุ Functional \& Non-Functional Requirements
      \item จัดทำ Use Cases \& Sequence Diagram
      \item วางโครงสร้าง System Architecture
      \item วางโครงสร้าง AI Workflow \& Prompt Architecture
      \item ออกแบบ Database
      \item ออกแบบ API Design
      \item ออกแบบ UI/UX Design (User Flow \& Wireframe)
      \item System prototype (PoC)
    \end{itemize}

  \item \textbf{M0: Foundation}
    \begin{itemize}
      \item ระบุกลุ่มเป้าหมาย ขอบเขต และ MVP
      \item กำหนด data model (workspace, project, story, variant, session, task, feedback, round, report)
      \item เลือก architecture: auth, LLM provider, prototype rendering, hosting
      \item สำรวจและเลือก design system และ UI kit
      \item กำหนด requirements ด้าน PDPA และข้อความขอ consent
      \item ตั้งค่า CI/CD, environments และ event analytics
      % \item การกำหนดดัชนีชี้วัดผล และแผนการดำเนินงาน
    \end{itemize}

  \item \textbf{M1: Secure Access}
    \begin{itemize}
      \item จัดทำ Landing page
      \item จัดทำ Register page
      \item การยืนยันตัวตนผ่านอีเมล
      \item การเข้าสู่ระบบพร้อมระบบจัดการข้อผิดพลาด
      \item ระบบรีเซ็ตรหัสผ่าน
      \item จัดทำระบบ routing ตามบทบาทของผู้ใช้
      \item ระบบควบคุมการเข้าถึง
    \end{itemize}

  \item \textbf{M2: Project Setup}
    \begin{itemize}
      \item ระบบสร้าง workspace
      \item ระบบสร้างโปรเจกต์
      \item ฟอร์มกำหนดความต้องการทางการออกแบบ
      \item ระบบจัดเก็บความต้องการทางการออกแบบเพื่อใช้เป็น reusable AI context
    \end{itemize}

  \item \textbf{M3: AI Design Generation}
    \begin{itemize}
      \item การสร้าง User story และ Problem statement
      \item เครื่องมือสร้างคำสั่งที่รวมความต้องการทางการออกแบบและ user story
      \item การสร้าง two-alternative UI
      \item การปรับปรุง UI (Chat-based)
      \item การดูตัวอย่างรูปแบบหน้าจอ
      \item การอนุมัติและมอบหมายงานให้ผู้เข้าร่วมทดสอบ
      \item การเริ่มเซสชันการทดสอบ
      % \item ชุดข้อมูลสำหรับการประเมินคุณภาพผลลัพธ์ของ AI
    \end{itemize}

  \item \textbf{M4: Testing Space}
    \begin{itemize}
      \item การเข้าสู่เซสชันด้วยรหัสหรือลิงก์ พร้อมระบบตรวจสอบความถูกต้องของอายุการใช้งาน
      \item สถานะสิ้นสุดของระบบ
      \item การขอความยินยอมตามกฎหมายคุ้มครองข้อมูลส่วนบุคคล (PDPA) พร้อมระบบบันทึกประวัติการให้ความยินยอม
      \item การกำหนดผู้ใช้งานให้เข้าร่วมการทดสอบ
      \item การแนะนำการใช้งานเครื่องมือทดสอบสำหรับผู้ใช้ใหม่
      \item ระบบทำภารกิจพร้อมการบันทึกพฤติกรรมการโต้ตอบของผู้ใช้
      \item การแสดงความคิดเห็น การระบุตำแหน่งคอมโพเนนต์ การแสดงความคิดเห็นทั้งหน้าจอ หรือการข้ามการแสดงความคิดเห็น
      \item วงจรการทำภารกิจที่เหลืออยู่จนกระทั่งสิ้นสุดการทดสอบของผู้ทดสอบ
      \item การทดสอบข้ามเบราว์เซอร์
    \end{itemize}

  \item \textbf{M5: Feedback}
    \begin{itemize}
      \item การติดตามเซสชันการทดสอบ
      \item การรวบรวมและประมวลผลฟีดแบ็ก
      \item การวิเคราะห์ด้วย AI การหาประเด็นหลัก การติดป้ายกำกับ และการสรุปผล
      \item หน้าจอเพื่อการตรวจสอบของนักออกแบบ การเปลี่ยนป้ายกำกับ การรวม การแยก การจดบันทึก และการทำเครื่องหมาย "ไม่แก้ไข" พร้อมระบุเหตุผล
      \item วงจรการอนุมัติหรือการปรับปรุงแก้ไข
      \item การนำฟีดแบ็กที่เลือกไปป้อนเป็นพรอมต์สำหรับรอบถัดไป
      \item การตัดสินใจสิ้นสุดเซสชัน การทำเวอร์ชันของรอบการทำงาน และประวัติการทำงาน
      \item การทดสอบวงจรการทำงานแบบต้นทางถึงปลายทาง
    \end{itemize}

  \item \textbf{M6: Stakeholder Report}
    \begin{itemize}
      \item รายการโปรเจกต์ (แบบอ่านอย่างเดียว สำหรับงานที่อนุมัติแล้วเท่านั้น)
      \item มุมมองผู้ใช้งานที่แตกต่างกัน
      \item การส่งมอบงานให้นักพัฒนา
      \item สถานะสำหรับผู้จัดการโครงการ
      \item เครื่องมือสร้างรายงาน
      \item การจัดทำรายงานเฉพาะสำหรับแต่ละบทบาท
      \item การส่งออกและการแชร์รายงาน
    \end{itemize}

  \item \textbf{M7: Pilot and Launch}
    \begin{itemize}
      \item ทดสอบ E2E และ UAT
      \item ทดสอบ Security test และ PDPA audit
      \item ทดสอบ Load test ด้วยผู้ทดสอบพร้อมกัน
      \item ทดลองใช้งานกับทีมจริง 2–3 ทีมและเก็บข้อเสนอแนะ
      \item จัดทำคู่มือการใช้งาน การเริ่มต้นใช้งาน และกระบวนการสนับสนุน
      \item ติดตามระบบและแดชบอร์ดตัวชี้วัดทางธุรกิจ
      \item Go live
    \end{itemize}
\end{enumerate}

\subsection{Gantt Chart}
% !TeX root = ../../progress-sem1.tex
% Generated by tools/gantt_to_latex.py from the Gantt sheet. Edit the sheet and regenerate,
% rather than editing this file by hand.
\begingroup
\footnotesize
\setlength{\tabcolsep}{2pt}
\renewcommand{\arraystretch}{1.15}
\begin{longtable}{|>{\RaggedRight\arraybackslash}p{4.4cm}|@{}>{\centering\arraybackslash}p{\dimexpr(\linewidth-4.4cm-1.9cm-4\tabcolsep-30\arrayrulewidth)/27\relax}@{}|@{}>{\centering\arraybackslash}p{\dimexpr(\linewidth-4.4cm-1.9cm-4\tabcolsep-30\arrayrulewidth)/27\relax}@{}|@{}>{\centering\arraybackslash}p{\dimexpr(\linewidth-4.4cm-1.9cm-4\tabcolsep-30\arrayrulewidth)/27\relax}@{}|@{}>{\centering\arraybackslash}p{\dimexpr(\linewidth-4.4cm-1.9cm-4\tabcolsep-30\arrayrulewidth)/27\relax}@{}|@{}>{\centering\arraybackslash}p{\dimexpr(\linewidth-4.4cm-1.9cm-4\tabcolsep-30\arrayrulewidth)/27\relax}@{}|@{}>{\centering\arraybackslash}p{\dimexpr(\linewidth-4.4cm-1.9cm-4\tabcolsep-30\arrayrulewidth)/27\relax}@{}|@{}>{\centering\arraybackslash}p{\dimexpr(\linewidth-4.4cm-1.9cm-4\tabcolsep-30\arrayrulewidth)/27\relax}@{}|@{}>{\centering\arraybackslash}p{\dimexpr(\linewidth-4.4cm-1.9cm-4\tabcolsep-30\arrayrulewidth)/27\relax}@{}|@{}>{\centering\arraybackslash}p{\dimexpr(\linewidth-4.4cm-1.9cm-4\tabcolsep-30\arrayrulewidth)/27\relax}@{}|@{}>{\centering\arraybackslash}p{\dimexpr(\linewidth-4.4cm-1.9cm-4\tabcolsep-30\arrayrulewidth)/27\relax}@{}|@{}>{\centering\arraybackslash}p{\dimexpr(\linewidth-4.4cm-1.9cm-4\tabcolsep-30\arrayrulewidth)/27\relax}@{}|@{}>{\centering\arraybackslash}p{\dimexpr(\linewidth-4.4cm-1.9cm-4\tabcolsep-30\arrayrulewidth)/27\relax}@{}|@{}>{\centering\arraybackslash}p{\dimexpr(\linewidth-4.4cm-1.9cm-4\tabcolsep-30\arrayrulewidth)/27\relax}@{}|@{}>{\centering\arraybackslash}p{\dimexpr(\linewidth-4.4cm-1.9cm-4\tabcolsep-30\arrayrulewidth)/27\relax}@{}|@{}>{\centering\arraybackslash}p{\dimexpr(\linewidth-4.4cm-1.9cm-4\tabcolsep-30\arrayrulewidth)/27\relax}@{}|@{}>{\centering\arraybackslash}p{\dimexpr(\linewidth-4.4cm-1.9cm-4\tabcolsep-30\arrayrulewidth)/27\relax}@{}|@{}>{\centering\arraybackslash}p{\dimexpr(\linewidth-4.4cm-1.9cm-4\tabcolsep-30\arrayrulewidth)/27\relax}@{}|@{}>{\centering\arraybackslash}p{\dimexpr(\linewidth-4.4cm-1.9cm-4\tabcolsep-30\arrayrulewidth)/27\relax}@{}|@{}>{\centering\arraybackslash}p{\dimexpr(\linewidth-4.4cm-1.9cm-4\tabcolsep-30\arrayrulewidth)/27\relax}@{}|@{}>{\centering\arraybackslash}p{\dimexpr(\linewidth-4.4cm-1.9cm-4\tabcolsep-30\arrayrulewidth)/27\relax}@{}|@{}>{\centering\arraybackslash}p{\dimexpr(\linewidth-4.4cm-1.9cm-4\tabcolsep-30\arrayrulewidth)/27\relax}@{}|@{}>{\centering\arraybackslash}p{\dimexpr(\linewidth-4.4cm-1.9cm-4\tabcolsep-30\arrayrulewidth)/27\relax}@{}|@{}>{\centering\arraybackslash}p{\dimexpr(\linewidth-4.4cm-1.9cm-4\tabcolsep-30\arrayrulewidth)/27\relax}@{}|@{}>{\centering\arraybackslash}p{\dimexpr(\linewidth-4.4cm-1.9cm-4\tabcolsep-30\arrayrulewidth)/27\relax}@{}|@{}>{\centering\arraybackslash}p{\dimexpr(\linewidth-4.4cm-1.9cm-4\tabcolsep-30\arrayrulewidth)/27\relax}@{}|@{}>{\centering\arraybackslash}p{\dimexpr(\linewidth-4.4cm-1.9cm-4\tabcolsep-30\arrayrulewidth)/27\relax}@{}|@{}>{\centering\arraybackslash}p{\dimexpr(\linewidth-4.4cm-1.9cm-4\tabcolsep-30\arrayrulewidth)/27\relax}@{}|>{\RaggedRight\arraybackslash}p{1.9cm}|}
\caption{แผนการดำเนินงานในภาคการศึกษาที่ 1}\label{tab:gantt-semester-1}\\
\hline
\multirow{2}{*}{\textbf{รายละเอียดงาน}} & \multicolumn{5}{c|}{\textbf{ก.ค. 69}} & \multicolumn{4}{c|}{\textbf{ส.ค. 69}} & \multicolumn{5}{c|}{\textbf{ก.ย. 69}} & \multicolumn{4}{c|}{\textbf{ต.ค. 69}} & \multicolumn{4}{c|}{\textbf{พ.ย. 69}} & \multicolumn{5}{c|}{\textbf{ธ.ค. 69}} & \multirow{2}{*}{\textbf{ผู้รับผิดชอบ}} \\
\cline{2-28}
 & 1 & 2 & 3 & 4 & 5 & 1 & 2 & 3 & 4 & 1 & 2 & 3 & 4 & 5 & 1 & 2 & 3 & 4 & 1 & 2 & 3 & 4 & 1 & 2 & 3 & 4 & 5 & \\
\hline
\endfirsthead
\multicolumn{29}{l}{\textit{ตารางที่ \thetable\ (ต่อ)}}\\
\hline
\multirow{2}{*}{\textbf{รายละเอียดงาน}} & \multicolumn{5}{c|}{\textbf{ก.ค. 69}} & \multicolumn{4}{c|}{\textbf{ส.ค. 69}} & \multicolumn{5}{c|}{\textbf{ก.ย. 69}} & \multicolumn{4}{c|}{\textbf{ต.ค. 69}} & \multicolumn{4}{c|}{\textbf{พ.ย. 69}} & \multicolumn{5}{c|}{\textbf{ธ.ค. 69}} & \multirow{2}{*}{\textbf{ผู้รับผิดชอบ}} \\
\cline{2-28}
 & 1 & 2 & 3 & 4 & 5 & 1 & 2 & 3 & 4 & 1 & 2 & 3 & 4 & 5 & 1 & 2 & 3 & 4 & 1 & 2 & 3 & 4 & 1 & 2 & 3 & 4 & 5 & \\
\hline
\endhead
\multicolumn{29}{|l|}{\textbf{Documentation}} \\ \hline
ศึกษาค้นคว้าหัวข้อโครงงาน & \cellcolor{ganttbar} & \cellcolor{ganttbar} & \cellcolor{ganttbar} & \cellcolor{ganttbar} & \cellcolor{ganttbar} & \cellcolor{ganttbar} & \cellcolor{ganttbar} &  &  &  &  &  &  &  &  &  &  &  &  &  &  &  &  &  &  &  &  & ทุกคน \\ \hline
จัดทำเอกสาร Project Idea &  &  &  &  &  & \cellcolor{ganttbar} & \cellcolor{ganttbar} & \cellcolor{ganttbar} & \cellcolor{ganttbar} &  &  &  &  &  &  &  &  &  &  &  &  &  &  &  &  &  &  & ทุกคน \\ \hline
เอกสารและงานวิจัยที่เกี่ยวข้อง &  &  &  &  &  &  &  &  &  & \cellcolor{ganttbar} & \cellcolor{ganttbar} & \cellcolor{ganttbar} &  &  &  &  &  &  &  &  &  &  &  &  &  &  &  & ภูบดินทร์, นาวิน \\ \hline
จัดทำเอกสาร Proposal report &  &  &  &  &  &  &  &  &  & \cellcolor{ganttbar} & \cellcolor{ganttbar} & \cellcolor{ganttbar} & \cellcolor{ganttbar} & \cellcolor{ganttbar} &  &  &  &  &  &  &  &  &  &  &  &  &  & ทุกคน \\ \hline
นำเสนอ Proposal &  &  &  &  &  &  &  &  &  &  &  &  &  & \cellcolor{ganttbar} & \cellcolor{ganttbar} &  &  &  &  &  &  &  &  &  &  &  &  & ทุกคน \\ \hline
จัดทำเอกสาร Progress report ภาคเรียนที่ 1 &  &  &  &  &  &  &  &  &  &  &  &  &  &  &  &  &  &  &  & \cellcolor{ganttbar} & \cellcolor{ganttbar} & \cellcolor{ganttbar} & \cellcolor{ganttbar} &  &  &  &  & ทุกคน \\ \hline
นำเสนอ Progress ภาคเรียนที่ 1 &  &  &  &  &  &  &  &  &  &  &  &  &  &  &  &  &  &  &  &  &  &  & \cellcolor{ganttbar} & \cellcolor{ganttbar} &  &  &  & ทุกคน \\ \hline
\multicolumn{29}{|l|}{\textbf{Requirements Analysis \& System Design}} \\ \hline
สัมภาษณ์ผู้มีส่วนได้ส่วนเสีย &  &  &  &  &  &  &  &  &  & \cellcolor{ganttbar} & \cellcolor{ganttbar} &  &  &  &  &  &  &  &  &  &  &  &  &  &  &  &  & กัญญ์ญาพัชร, ปัณณวิชญ์ \\ \hline
ระบุ Requirements &  &  &  &  &  &  &  &  &  &  & \cellcolor{ganttbar} & \cellcolor{ganttbar} &  &  &  &  &  &  &  &  &  &  &  &  &  &  &  & ปัณณวิชญ์, กัญญ์ญาพัชร \\ \hline
ระบุ Functional \& Non-Functional Requirements &  &  &  &  &  &  &  &  &  &  & \cellcolor{ganttbar} & \cellcolor{ganttbar} &  &  &  &  &  &  &  &  &  &  &  &  &  &  &  & กัญญ์ญาพัชร, ปัณณวิชญ์ \\ \hline
จัดทำ Use Cases \& Sequence Diagram &  &  &  &  &  &  &  &  &  &  &  &  & \cellcolor{ganttbar} &  &  &  &  &  &  &  &  &  &  &  &  &  &  & นาวิน, กัญญ์ญาพัชร \\ \hline
วางโครงสร้าง System Architecture &  &  &  &  &  &  &  &  &  &  &  & \cellcolor{ganttbar} & \cellcolor{ganttbar} &  &  &  &  &  &  &  &  &  &  &  &  &  &  & ภูบดินทร์ \\ \hline
วางโครงสร้าง AI Workflow \& Prompt Architecture &  &  &  &  &  &  &  &  &  &  &  &  &  &  &  & \cellcolor{ganttbar} & \cellcolor{ganttbar} &  &  &  &  &  &  &  &  &  &  & ภูบดินทร์, ปัณณวิชญ์ \\ \hline
ออกแบบ Database &  &  &  &  &  &  &  &  &  &  &  &  &  &  &  & \cellcolor{ganttbar} & \cellcolor{ganttbar} &  &  &  &  &  &  &  &  &  &  & ภูบดินทร์, นาวิน \\ \hline
ออกแบบ API Design &  &  &  &  &  &  &  &  &  &  &  &  &  &  &  &  & \cellcolor{ganttbar} & \cellcolor{ganttbar} &  &  &  &  &  &  &  &  &  & ภูบดินทร์ \\ \hline
ออกแบบ UI/UX Design (User Flow \& Wireframe) &  &  &  &  &  &  &  &  &  &  &  & \cellcolor{ganttbar} &  &  &  &  &  &  &  &  &  &  &  &  &  &  &  & ปัณณวิชญ์ \\ \hline
System prototype (PoC) &  &  &  &  &  &  &  &  &  &  &  &  &  &  &  &  & \cellcolor{ganttbar} & \cellcolor{ganttbar} &  &  &  &  &  &  &  &  &  & ทุกคน \\ \hline
\multicolumn{29}{|l|}{\textbf{M0: Foundation}} \\ \hline
ระบุกลุ่มเป้าหมาย ขอบเขต และ MVP &  &  &  &  &  &  &  &  &  &  & \cellcolor{ganttbar} & \cellcolor{ganttbar} &  &  &  &  &  &  &  &  &  &  &  &  &  &  &  & กัญญ์ญาพัชร, ปัณณวิชญ์ \\ \hline
กำหนด data model (workspace, project, story, variant, session, task, feedback, round, report) &  &  &  &  &  &  &  &  &  &  &  &  &  &  &  & \cellcolor{ganttbar} & \cellcolor{ganttbar} &  &  &  &  &  &  &  &  &  &  & นาวิน, ภูบดินทร์ \\ \hline
เลือก architecture: auth, LLM provider, prototype rendering, hosting &  &  &  &  &  &  &  &  &  &  &  &  &  &  &  & \cellcolor{ganttbar} & \cellcolor{ganttbar} &  &  &  &  &  &  &  &  &  &  & ภูบดินทร์, นาวิน \\ \hline
สำรวจและเลือก design system และ UI kit &  &  &  &  &  &  &  &  &  &  &  &  &  &  &  &  & \cellcolor{ganttbar} & \cellcolor{ganttbar} &  &  &  &  &  &  &  &  &  & ปัณณวิชญ์ \\ \hline
กำหนด requirements ด้าน PDPA และข้อความขอ consent &  &  &  &  &  &  &  &  &  &  &  &  &  &  &  & \cellcolor{ganttbar} & \cellcolor{ganttbar} &  &  &  &  &  &  &  &  &  &  & กัญญ์ญาพัชร \\ \hline
ตั้งค่า CI/CD, environments และ event analytics &  &  &  &  &  &  &  &  &  &  &  &  &  &  &  &  & \cellcolor{ganttbar} & \cellcolor{ganttbar} &  &  &  &  &  &  &  &  &  & ภูบดินทร์ \\ \hline
การกำหนดดัชนีชี้วัดผล และแผนการดำเนินงาน &  &  &  &  &  &  &  &  &  &  &  &  &  &  &  & \cellcolor{ganttbar} &  &  &  &  &  &  &  &  &  &  &  & ทุกคน \\ \hline
\multicolumn{29}{|l|}{\textbf{M1: Secure Access}} \\ \hline
จัดทำ Landing page &  &  &  &  &  &  &  &  &  &  &  &  &  &  &  &  & \cellcolor{ganttbar} & \cellcolor{ganttbar} &  &  &  &  &  &  &  &  &  & ปัณณวิชญ์, กัญญ์ญาพัชร, นาวิน \\ \hline
จัดทำ Register page &  &  &  &  &  &  &  &  &  &  &  &  &  &  &  &  & \cellcolor{ganttbar} & \cellcolor{ganttbar} & \cellcolor{ganttbar} &  &  &  &  &  &  &  &  & กัญญ์ญาพัชร, นาวิน, ภูบดินทร์ \\ \hline
การยืนยันตัวตนผ่านอีเมล &  &  &  &  &  &  &  &  &  &  &  &  &  &  &  &  &  &  & \cellcolor{ganttbar} & \cellcolor{ganttbar} &  &  &  &  &  &  &  & ภูบดินทร์ \\ \hline
การเข้าสู่ระบบพร้อมระบบจัดการข้อผิดพลาด &  &  &  &  &  &  &  &  &  &  &  &  &  &  &  &  &  & \cellcolor{ganttbar} & \cellcolor{ganttbar} & \cellcolor{ganttbar} &  &  &  &  &  &  &  & กัญญ์ญาพัชร, นาวิน, ภูบดินทร์ \\ \hline
ระบบรีเซ็ตรหัสผ่าน &  &  &  &  &  &  &  &  &  &  &  &  &  &  &  &  &  &  &  & \cellcolor{ganttbar} & \cellcolor{ganttbar} & \cellcolor{ganttbar} &  &  &  &  &  & กัญญ์ญาพัชร, นาวิน, ภูบดินทร์ \\ \hline
ระบบควบคุมการเข้าถึง &  &  &  &  &  &  &  &  &  &  &  &  &  &  &  &  &  &  &  & \cellcolor{ganttbar} & \cellcolor{ganttbar} & \cellcolor{ganttbar} &  &  &  &  &  & ภูบดินทร์ \\ \hline
จัดทำระบบ routing ตามบทบาทของผู้ใช้ &  &  &  &  &  &  &  &  &  &  &  &  &  &  &  &  &  &  &  &  &  & \cellcolor{ganttbar} & \cellcolor{ganttbar} & \cellcolor{ganttbar} &  &  &  & กัญญ์ญาพัชร, นาวิน \\ \hline
ทดสอบระบบยืนยันตัวตนและการควบคุมการเข้าถึง &  &  &  &  &  &  &  &  &  &  &  &  &  &  &  &  &  &  &  &  &  & \cellcolor{ganttbar} & \cellcolor{ganttbar} & \cellcolor{ganttbar} &  &  &  & ภูบดินทร์ \\ \hline
\multicolumn{29}{|l|}{\textbf{M2: Project Setup}} \\ \hline
ระบบสร้าง workspace &  &  &  &  &  &  &  &  &  &  &  &  &  &  &  &  &  &  &  &  &  & \cellcolor{ganttbar} & \cellcolor{ganttbar} & \cellcolor{ganttbar} &  &  &  & กัญญ์ญาพัชร, นาวิน, ภูบดินทร์ \\ \hline
ระบบสร้างโปรเจกต์ &  &  &  &  &  &  &  &  &  &  &  &  &  &  &  &  &  &  &  &  &  &  & \cellcolor{ganttbar} & \cellcolor{ganttbar} &  &  &  & กัญญ์ญาพัชร, นาวิน, ภูบดินทร์ \\ \hline
ฟอร์มกำหนดความต้องการทางการออกแบบ &  &  &  &  &  &  &  &  &  &  &  &  &  &  &  &  &  &  &  &  &  &  & \cellcolor{ganttbar} & \cellcolor{ganttbar} & \cellcolor{ganttbar} & \cellcolor{ganttbar} &  & ปัณณวิชญ์, กัญญ์ญาพัชร, นาวิน \\ \hline
ระบบจัดเก็บความต้องการทางการออกแบบเพื่อใช้เป็น reusable AI context &  &  &  &  &  &  &  &  &  &  &  &  &  &  &  &  &  &  &  &  &  &  &  & \cellcolor{ganttbar} & \cellcolor{ganttbar} & \cellcolor{ganttbar} &  & ภูบดินทร์ \\ \hline
\multicolumn{29}{|l|}{\textbf{M3: AI Design Generation}} \\ \hline
การสร้าง User story และ Problem statement &  &  &  &  &  &  &  &  &  &  &  &  &  &  &  &  &  &  &  &  &  &  &  & \cellcolor{ganttbar} & \cellcolor{ganttbar} & \cellcolor{ganttbar} &  & กัญญ์ญาพัชร, นาวิน, ภูบดินทร์ \\ \hline
เครื่องมือสร้างคำสั่งที่รวมความต้องการทางการออกแบบและ user story &  &  &  &  &  &  &  &  &  &  &  &  &  &  &  &  &  &  &  &  &  &  &  &  & \cellcolor{ganttbar} & \cellcolor{ganttbar} & \cellcolor{ganttbar} & ภูบดินทร์ \\ \hline
การสร้าง two-alternative UI &  &  &  &  &  &  &  &  &  &  &  &  &  &  &  &  &  &  &  &  &  &  &  &  &  & \cellcolor{ganttbar} & \cellcolor{ganttbar} & ภูบดินทร์ \\ \hline
การปรับปรุง UI (Chat-based) &  &  &  &  &  &  &  &  &  &  &  &  &  &  &  &  &  &  &  &  &  &  &  &  &  &  & \cellcolor{ganttbar} & กัญญ์ญาพัชร, นาวิน, ภูบดินทร์ \\ \hline
การดูตัวอย่างรูปแบบหน้าจอ &  &  &  &  &  &  &  &  &  &  &  &  &  &  &  &  &  &  &  &  &  &  &  &  &  & \cellcolor{ganttbar} & \cellcolor{ganttbar} & กัญญ์ญาพัชร, นาวิน \\ \hline
ชุดข้อมูลสำหรับการประเมินคุณภาพผลลัพธ์ของ AI &  &  &  &  &  &  &  &  &  &  &  &  &  &  &  &  &  &  &  &  &  &  &  &  &  &  & \cellcolor{ganttbar} & ปัณณวิชญ์, ภูบดินทร์ \\ \hline
\end{longtable}
\endgroup

\begingroup
\footnotesize
\setlength{\tabcolsep}{2pt}
\renewcommand{\arraystretch}{1.15}
\begin{longtable}{|>{\RaggedRight\arraybackslash}p{4.4cm}|@{}>{\centering\arraybackslash}p{\dimexpr(\linewidth-4.4cm-1.9cm-4\tabcolsep-24\arrayrulewidth)/21\relax}@{}|@{}>{\centering\arraybackslash}p{\dimexpr(\linewidth-4.4cm-1.9cm-4\tabcolsep-24\arrayrulewidth)/21\relax}@{}|@{}>{\centering\arraybackslash}p{\dimexpr(\linewidth-4.4cm-1.9cm-4\tabcolsep-24\arrayrulewidth)/21\relax}@{}|@{}>{\centering\arraybackslash}p{\dimexpr(\linewidth-4.4cm-1.9cm-4\tabcolsep-24\arrayrulewidth)/21\relax}@{}|@{}>{\centering\arraybackslash}p{\dimexpr(\linewidth-4.4cm-1.9cm-4\tabcolsep-24\arrayrulewidth)/21\relax}@{}|@{}>{\centering\arraybackslash}p{\dimexpr(\linewidth-4.4cm-1.9cm-4\tabcolsep-24\arrayrulewidth)/21\relax}@{}|@{}>{\centering\arraybackslash}p{\dimexpr(\linewidth-4.4cm-1.9cm-4\tabcolsep-24\arrayrulewidth)/21\relax}@{}|@{}>{\centering\arraybackslash}p{\dimexpr(\linewidth-4.4cm-1.9cm-4\tabcolsep-24\arrayrulewidth)/21\relax}@{}|@{}>{\centering\arraybackslash}p{\dimexpr(\linewidth-4.4cm-1.9cm-4\tabcolsep-24\arrayrulewidth)/21\relax}@{}|@{}>{\centering\arraybackslash}p{\dimexpr(\linewidth-4.4cm-1.9cm-4\tabcolsep-24\arrayrulewidth)/21\relax}@{}|@{}>{\centering\arraybackslash}p{\dimexpr(\linewidth-4.4cm-1.9cm-4\tabcolsep-24\arrayrulewidth)/21\relax}@{}|@{}>{\centering\arraybackslash}p{\dimexpr(\linewidth-4.4cm-1.9cm-4\tabcolsep-24\arrayrulewidth)/21\relax}@{}|@{}>{\centering\arraybackslash}p{\dimexpr(\linewidth-4.4cm-1.9cm-4\tabcolsep-24\arrayrulewidth)/21\relax}@{}|@{}>{\centering\arraybackslash}p{\dimexpr(\linewidth-4.4cm-1.9cm-4\tabcolsep-24\arrayrulewidth)/21\relax}@{}|@{}>{\centering\arraybackslash}p{\dimexpr(\linewidth-4.4cm-1.9cm-4\tabcolsep-24\arrayrulewidth)/21\relax}@{}|@{}>{\centering\arraybackslash}p{\dimexpr(\linewidth-4.4cm-1.9cm-4\tabcolsep-24\arrayrulewidth)/21\relax}@{}|@{}>{\centering\arraybackslash}p{\dimexpr(\linewidth-4.4cm-1.9cm-4\tabcolsep-24\arrayrulewidth)/21\relax}@{}|@{}>{\centering\arraybackslash}p{\dimexpr(\linewidth-4.4cm-1.9cm-4\tabcolsep-24\arrayrulewidth)/21\relax}@{}|@{}>{\centering\arraybackslash}p{\dimexpr(\linewidth-4.4cm-1.9cm-4\tabcolsep-24\arrayrulewidth)/21\relax}@{}|@{}>{\centering\arraybackslash}p{\dimexpr(\linewidth-4.4cm-1.9cm-4\tabcolsep-24\arrayrulewidth)/21\relax}@{}|@{}>{\centering\arraybackslash}p{\dimexpr(\linewidth-4.4cm-1.9cm-4\tabcolsep-24\arrayrulewidth)/21\relax}@{}|>{\RaggedRight\arraybackslash}p{1.9cm}|}
\caption{แผนการดำเนินงานในภาคการศึกษาที่ 2}\label{tab:gantt-semester-2}\\
\hline
\multirow{2}{*}{\textbf{รายละเอียดงาน}} & \multicolumn{4}{c|}{\textbf{ม.ค. 70}} & \multicolumn{4}{c|}{\textbf{ก.พ. 70}} & \multicolumn{5}{c|}{\textbf{มี.ค. 70}} & \multicolumn{4}{c|}{\textbf{เม.ย. 70}} & \multicolumn{4}{c|}{\textbf{พ.ค. 70}} & \multirow{2}{*}{\textbf{ผู้รับผิดชอบ}} \\
\cline{2-22}
 & 1 & 2 & 3 & 4 & 1 & 2 & 3 & 4 & 1 & 2 & 3 & 4 & 5 & 1 & 2 & 3 & 4 & 1 & 2 & 3 & 4 & \\
\hline
\endfirsthead
\multicolumn{23}{l}{\textit{ตารางที่ \thetable\ (ต่อ)}}\\
\hline
\multirow{2}{*}{\textbf{รายละเอียดงาน}} & \multicolumn{4}{c|}{\textbf{ม.ค. 70}} & \multicolumn{4}{c|}{\textbf{ก.พ. 70}} & \multicolumn{5}{c|}{\textbf{มี.ค. 70}} & \multicolumn{4}{c|}{\textbf{เม.ย. 70}} & \multicolumn{4}{c|}{\textbf{พ.ค. 70}} & \multirow{2}{*}{\textbf{ผู้รับผิดชอบ}} \\
\cline{2-22}
 & 1 & 2 & 3 & 4 & 1 & 2 & 3 & 4 & 1 & 2 & 3 & 4 & 5 & 1 & 2 & 3 & 4 & 1 & 2 & 3 & 4 & \\
\hline
\endhead
\multicolumn{23}{|l|}{\textbf{Documentation}} \\ \hline
จัดทำเอกสาร Progress report ภาคเรียนที่ 2 &  &  &  &  & \cellcolor{ganttbar} & \cellcolor{ganttbar} & \cellcolor{ganttbar} & \cellcolor{ganttbar} &  &  &  &  &  &  &  &  &  &  &  &  &  & ทุกคน \\ \hline
นำเสนอ Progress ภาคเรียนที่ 2 &  &  &  &  &  &  &  &  & \cellcolor{ganttbar} &  &  &  &  &  &  &  &  &  &  &  &  & ทุกคน \\ \hline
จัดทำเอกสาร Final Report &  &  &  &  &  &  &  &  &  &  &  &  & \cellcolor{ganttbar} & \cellcolor{ganttbar} & \cellcolor{ganttbar} & \cellcolor{ganttbar} & \cellcolor{ganttbar} &  &  &  &  & ทุกคน \\ \hline
นำเสนอ Final Defense &  &  &  &  &  &  &  &  &  &  &  &  & \cellcolor{ganttbar} & \cellcolor{ganttbar} & \cellcolor{ganttbar} & \cellcolor{ganttbar} & \cellcolor{ganttbar} &  &  &  &  & ทุกคน \\ \hline
\multicolumn{23}{|l|}{\textbf{M3: AI Design Generation}} \\ \hline
เครื่องมือสร้างคำสั่งที่รวมความต้องการทางการออกแบบและ user story & \cellcolor{ganttbar} &  &  &  &  &  &  &  &  &  &  &  &  &  &  &  &  &  &  &  &  & ภูบดินทร์ \\ \hline
การสร้าง two-alternative UI & \cellcolor{ganttbar} & \cellcolor{ganttbar} &  &  &  &  &  &  &  &  &  &  &  &  &  &  &  &  &  &  &  & ภูบดินทร์ \\ \hline
การปรับปรุง UI (Chat-based) & \cellcolor{ganttbar} & \cellcolor{ganttbar} &  &  &  &  &  &  &  &  &  &  &  &  &  &  &  &  &  &  &  & กัญญ์ญาพัชร, นาวิน, ภูบดินทร์ \\ \hline
การดูตัวอย่างรูปแบบหน้าจอ & \cellcolor{ganttbar} & \cellcolor{ganttbar} &  &  &  &  &  &  &  &  &  &  &  &  &  &  &  &  &  &  &  & กัญญ์ญาพัชร, นาวิน \\ \hline
การอนุมัติและมอบหมายงานให้ผู้เข้าร่วมทดสอบ &  & \cellcolor{ganttbar} & \cellcolor{ganttbar} &  &  &  &  &  &  &  &  &  &  &  &  &  &  &  &  &  &  & กัญญ์ญาพัชร, นาวิน, ภูบดินทร์ \\ \hline
การเริ่มเซสชันการทดสอบ &  & \cellcolor{ganttbar} & \cellcolor{ganttbar} &  &  &  &  &  &  &  &  &  &  &  &  &  &  &  &  &  &  & กัญญ์ญาพัชร, นาวิน, ภูบดินทร์ \\ \hline
ชุดข้อมูลสำหรับการประเมินคุณภาพผลลัพธ์ของ AI & \cellcolor{ganttbar} & \cellcolor{ganttbar} & \cellcolor{ganttbar} &  &  &  &  &  &  &  &  &  &  &  &  &  &  &  &  &  &  & ปัณณวิชญ์, ภูบดินทร์ \\ \hline
\multicolumn{23}{|l|}{\textbf{M4: Testing Space}} \\ \hline
การเข้าสู่เซสชันด้วยรหัสหรือลิงก์ พร้อมระบบตรวจสอบความถูกต้องของอายุการใช้งาน &  &  &  &  & \cellcolor{ganttbar} & \cellcolor{ganttbar} &  &  &  &  &  &  &  &  &  &  &  &  &  &  &  & กัญญ์ญาพัชร, นาวิน, ภูบดินทร์ \\ \hline
สถานะสิ้นสุดของระบบ &  &  &  &  &  & \cellcolor{ganttbar} & \cellcolor{ganttbar} &  &  &  &  &  &  &  &  &  &  &  &  &  &  & กัญญ์ญาพัชร, นาวิน \\ \hline
การขอความยินยอมตามกฎหมายคุ้มครองข้อมูลส่วนบุคคล (PDPA) พร้อมระบบบันทึกประวัติการให้ความยินยอม &  &  &  &  &  &  & \cellcolor{ganttbar} & \cellcolor{ganttbar} &  &  &  &  &  &  &  &  &  &  &  &  &  & กัญญ์ญาพัชร, นาวิน, ภูบดินทร์ \\ \hline
การกำหนดผู้ใช้งานให้เข้าร่วมการทดสอบ &  &  &  &  &  &  &  & \cellcolor{ganttbar} & \cellcolor{ganttbar} &  &  &  &  &  &  &  &  &  &  &  &  & ภูบดินทร์ \\ \hline
การแนะนำการใช้งานเครื่องมือทดสอบสำหรับผู้ใช้ใหม่ &  &  &  &  &  &  & \cellcolor{ganttbar} & \cellcolor{ganttbar} &  &  &  &  &  &  &  &  &  &  &  &  &  & ปัณณวิชญ์, กัญญ์ญาพัชร, นาวิน \\ \hline
ระบบทำภารกิจพร้อมการบันทึกพฤติกรรมการโต้ตอบของผู้ใช้ &  &  &  &  &  &  & \cellcolor{ganttbar} & \cellcolor{ganttbar} &  &  &  &  &  &  &  &  &  &  &  &  &  & กัญญ์ญาพัชร, นาวิน, ภูบดินทร์ \\ \hline
การทบทวนหน้าจอหลังทำภารกิจและคำถามเฉพาะจุด &  &  &  &  &  &  &  & \cellcolor{ganttbar} & \cellcolor{ganttbar} & \cellcolor{ganttbar} &  &  &  &  &  &  &  &  &  &  &  & กัญญ์ญาพัชร, นาวิน \\ \hline
การแสดงความคิดเห็น การระบุตำแหน่งคอมโพเนนต์ การแสดงความคิดเห็นทั้งหน้าจอ หรือการข้ามการแสดงความคิดเห็น &  &  &  &  &  &  &  & \cellcolor{ganttbar} & \cellcolor{ganttbar} & \cellcolor{ganttbar} & \cellcolor{ganttbar} &  &  &  &  &  &  &  &  &  &  & กัญญ์ญาพัชร, นาวิน \\ \hline
วงจรการทำภารกิจที่เหลืออยู่จนกระทั่งสิ้นสุดการทดสอบของผู้ทดสอบ &  &  &  &  &  &  &  &  & \cellcolor{ganttbar} & \cellcolor{ganttbar} &  &  &  &  &  &  &  &  &  &  &  & กัญญ์ญาพัชร, นาวิน \\ \hline
การทดสอบข้ามเบราว์เซอร์ &  &  &  &  &  &  &  &  &  & \cellcolor{ganttbar} & \cellcolor{ganttbar} & \cellcolor{ganttbar} &  &  &  &  &  &  &  &  &  & กัญญ์ญาพัชร, นาวิน \\ \hline
\multicolumn{23}{|l|}{\textbf{M5: Feedback}} \\ \hline
การติดตามเซสชันการทดสอบ &  &  &  &  &  &  &  &  &  & \cellcolor{ganttbar} & \cellcolor{ganttbar} & \cellcolor{ganttbar} &  &  &  &  &  &  &  &  &  & กัญญ์ญาพัชร, นาวิน \\ \hline
การรวบรวมและประมวลผลฟีดแบ็ก &  &  &  &  &  &  &  &  &  & \cellcolor{ganttbar} & \cellcolor{ganttbar} & \cellcolor{ganttbar} &  &  &  &  &  &  &  &  &  & ภูบดินทร์ \\ \hline
การวิเคราะห์ด้วย AI การหาประเด็นหลัก การติดป้ายกำกับ และการสรุปผล &  &  &  &  &  &  &  &  &  &  & \cellcolor{ganttbar} & \cellcolor{ganttbar} & \cellcolor{ganttbar} & \cellcolor{ganttbar} &  &  &  &  &  &  &  & ภูบดินทร์ \\ \hline
หน้าจอเพื่อการตรวจสอบของนักออกแบบ การเปลี่ยนป้ายกำกับ การรวม การแยก การจดบันทึก และการทำเครื่องหมาย "ไม่แก้ไข" พร้อมระบุเหตุผล &  &  &  &  &  &  &  &  &  &  & \cellcolor{ganttbar} & \cellcolor{ganttbar} & \cellcolor{ganttbar} & \cellcolor{ganttbar} &  &  &  &  &  &  &  & ปัณณวิชญ์, กัญญ์ญาพัชร, นาวิน \\ \hline
วงจรการอนุมัติหรือการปรับปรุงแก้ไข &  &  &  &  &  &  &  &  &  &  &  & \cellcolor{ganttbar} & \cellcolor{ganttbar} &  &  &  &  &  &  &  &  & กัญญ์ญาพัชร, นาวิน, ภูบดินทร์ \\ \hline
การนำฟีดแบ็กที่เลือกไปป้อนเป็นพรอมต์สำหรับรอบถัดไป &  &  &  &  &  &  &  &  &  &  &  & \cellcolor{ganttbar} & \cellcolor{ganttbar} & \cellcolor{ganttbar} &  &  &  &  &  &  &  & ภูบดินทร์ \\ \hline
การตัดสินใจสิ้นสุดเซสชัน การทำเวอร์ชันของรอบการทำงาน และประวัติการทำงาน &  &  &  &  &  &  &  &  &  & \cellcolor{ganttbar} & \cellcolor{ganttbar} & \cellcolor{ganttbar} &  &  &  &  &  &  &  &  &  & ภูบดินทร์ \\ \hline
การทดสอบวงจรการทำงานแบบต้นทางถึงปลายทาง &  &  &  &  &  &  &  &  &  &  &  & \cellcolor{ganttbar} & \cellcolor{ganttbar} &  &  &  &  &  &  &  &  & ทุกคน \\ \hline
\multicolumn{23}{|l|}{\textbf{M6: Stakeholder Report}} \\ \hline
รายการโปรเจกต์ (แบบอ่านอย่างเดียว สำหรับงานที่อนุมัติแล้วเท่านั้น) &  &  &  &  &  &  &  &  &  &  &  & \cellcolor{ganttbar} & \cellcolor{ganttbar} &  &  &  &  &  &  &  &  & กัญญ์ญาพัชร, นาวิน \\ \hline
มุมมองผู้ใช้งานที่แตกต่างกัน &  &  &  &  &  &  &  &  &  &  &  &  &  & \cellcolor{ganttbar} & \cellcolor{ganttbar} & \cellcolor{ganttbar} &  &  &  &  &  & ปัณณวิชญ์, กัญญ์ญาพัชร, นาวิน \\ \hline
การส่งมอบงานให้นักพัฒนา &  &  &  &  &  &  &  &  &  &  &  &  & \cellcolor{ganttbar} & \cellcolor{ganttbar} & \cellcolor{ganttbar} & \cellcolor{ganttbar} &  &  &  &  &  & กัญญ์ญาพัชร, นาวิน, ภูบดินทร์ \\ \hline
สถานะสำหรับผู้จัดการโครงการ &  &  &  &  &  &  &  &  &  &  &  &  & \cellcolor{ganttbar} & \cellcolor{ganttbar} & \cellcolor{ganttbar} &  &  &  &  &  &  & กัญญ์ญาพัชร, นาวิน \\ \hline
เครื่องมือสร้างรายงาน &  &  &  &  &  &  &  &  &  &  &  &  & \cellcolor{ganttbar} & \cellcolor{ganttbar} & \cellcolor{ganttbar} & \cellcolor{ganttbar} &  &  &  &  &  & ภูบดินทร์ \\ \hline
การจัดทำรายงานเฉพาะสำหรับแต่ละบทบาท &  &  &  &  &  &  &  &  &  &  &  &  &  & \cellcolor{ganttbar} & \cellcolor{ganttbar} & \cellcolor{ganttbar} & \cellcolor{ganttbar} &  &  &  &  & ปัณณวิชญ์, กัญญ์ญาพัชร, นาวิน \\ \hline
การส่งออกและการแชร์รายงาน &  &  &  &  &  &  &  &  &  &  &  &  &  &  &  & \cellcolor{ganttbar} & \cellcolor{ganttbar} &  &  &  &  & ภูบดินทร์ \\ \hline
\multicolumn{23}{|l|}{\textbf{M7: Pilot and Launch}} \\ \hline
ทดสอบ E2E และ UAT &  &  &  &  &  &  &  &  &  &  &  &  &  &  &  &  & \cellcolor{ganttbar} & \cellcolor{ganttbar} &  &  &  & ทุกคน \\ \hline
ทดสอบ Security test และ PDPA audit &  &  &  &  &  &  &  &  &  &  &  &  &  &  &  &  &  & \cellcolor{ganttbar} &  &  &  & ภูบดินทร์ \\ \hline
ทดสอบ Load test ด้วยผู้ทดสอบพร้อมกัน &  &  &  &  &  &  &  &  &  &  &  &  &  &  &  &  &  & \cellcolor{ganttbar} &  &  &  & ภูบดินทร์ \\ \hline
ทดลองใช้งานกับทีมจริง 2–3 ทีมและเก็บข้อเสนอแนะ &  &  &  &  &  &  &  &  &  &  &  &  &  &  &  &  &  & \cellcolor{ganttbar} & \cellcolor{ganttbar} & \cellcolor{ganttbar} & \cellcolor{ganttbar} & ทุกคน \\ \hline
จัดทำคู่มือการใช้งาน การเริ่มต้นใช้งาน และกระบวนการสนับสนุน &  &  &  &  &  &  &  &  &  &  &  &  &  &  &  &  & \cellcolor{ganttbar} & \cellcolor{ganttbar} &  &  &  & ปัณณวิชญ์ \\ \hline
ติดตามระบบและแดชบอร์ดตัวชี้วัดทางธุรกิจ &  &  &  &  &  &  &  &  &  &  &  &  &  &  &  &  & \cellcolor{ganttbar} & \cellcolor{ganttbar} &  &  &  & ภูบดินทร์ \\ \hline
Go live &  &  &  &  &  &  &  &  &  &  &  &  &  &  &  &  &  &  &  & \cellcolor{ganttbar} & \cellcolor{ganttbar} & ทุกคน \\ \hline
\end{longtable}
\endgroup


